\section{An empirical illustration}
\label{sec:empirical}

\begin{table}[t]
\centering
\footnotesize
\setlength{\tabcolsep}{3.5pt}
\caption{Twenty moving-average timing rules against buy-and-hold on the US market, in the 7 of 51 non-overlapping two-year windows in which at least one test rejects at 5\% (Null Zoo v2b, part G; $p$-values).}
\label{tab:empirical}
\begin{tabular}{lrrrrrrrr}
\toprule
 & \multicolumn{5}{c}{Mean block length 8} & \multicolumn{2}{c}{22} & \multicolumn{1}{c}{} \\
\cmidrule(lr){2-6}\cmidrule(lr){7-8}
Window & RC & SPA$_c$-U & SPA$_c$ & SPA$_c$-BT & SPA$_c$-NB & SPA$_c$-U & SPA$_c$ & LET \\
\midrule
Jul 1930--Apr 1932 & \textbf{0.038} & \textbf{0.038} & \textbf{0.020} & \textbf{0.030} & \textbf{0.027} & \textbf{0.015} & \textbf{0.008} & 0.148 \\
Dec 1958--Dec 1960 & 0.084 & 0.084 & 0.115 & 0.097 & 0.107 & \textbf{0.032} & 0.056 & 0.572 \\
Dec 1960--Dec 1962 & 0.072 & 0.072 & \textbf{0.047} & \textbf{0.050} & \textbf{0.024} & 0.101 & 0.084 & 0.486 \\
Jan 1969--Jan 1971 & \textbf{0.019} & \textbf{0.019} & \textbf{0.010} & \textbf{0.007} & \textbf{0.003} & \textbf{0.022} & \textbf{0.016} & \textbf{0.008} \\
Jan 1971--Jan 1973 & 0.063 & 0.063 & \textbf{0.049} & 0.115 & 0.094 & \textbf{0.031} & \textbf{0.032} & 0.190 \\
Jan 1973--Jan 1975 & \textbf{0.010} & \textbf{0.010} & \textbf{0.000} & \textbf{0.004} & \textbf{0.003} & \textbf{0.015} & \textbf{0.001} & \textbf{0.022} \\
Jan 1977--Jan 1979 & 0.078 & 0.078 & 0.074 & 0.068 & 0.055 & \textbf{0.042} & \textbf{0.048} & 0.248 \\
\midrule
Windows rejected & 3 & 3 & 5 & 4 & 4 & 6 & 5 & 2 \\
Full sample$^{\dagger}$ & \textbf{0.050} & \textbf{0.050} & \textbf{0.042} & \textbf{0.038} & \textbf{0.031} & \textbf{0.043} & \textbf{0.039} & 0.316 \\
\bottomrule
\end{tabular}
\par\smallskip\begin{minipage}{0.97\linewidth}\footnotesize $p$-values with 1{,}000 stationary-bootstrap resamples; bold: at most 0.05. LET is the \v{S}id\'ak-adjusted $t$ test on the best rule's Sharpe ratio, with no resampling. Floored tests count ties. A superscript gives the number of rules used when some rules never left the market in a window and so had no excess return. Windows rejected: of all 51 windows. $^{\dagger}$4 March 1927 to 31 August 2026, 26{,}117 days, mean block length 30 for the first five tests. Every window is in Table~\ref{tab:empirical-all}.\end{minipage}
\end{table}

\begin{table}[t]
\centering
\small
\caption{Published cross-sectional predictors declared superior to zero by Romano--Wolf stepwise tests at 5\%, in ten-year windows of monthly long-short returns (pre-registered illustration; counts of predictors).}
\label{tab:anomalies}
\begin{tabular}{lrrrr}
\toprule
Window & Fixed, $q=10$ & Fixed, $q=5$ & Bootstrap-$t$ & Reality Check \\
\midrule
1975--1984 ($k=162$) & 45 & 43 & 27 & 8 \\
1980--1989 ($k=179$) & 69 & 67 & 58 & 13 \\
1985--1994 ($k=188$) & 61 & 53 & 42 & 22 \\
1990--1999 ($k=193$) & 64 & 56 & 42 & 15 \\
1995--2004 ($k=202$) & 46 & 40 & 32 & 1 \\
2000--2009 ($k=206$) & 17 & 15 & 12 & 1 \\
2005--2014 ($k=206$) & 17 & 9 & 5 & 1 \\
2010--2019 ($k=207$) & 26 & 15 & 7 & 2 \\
2015--2024 ($k=195$) & 3 & 0 & 0 & 3 \\
\bottomrule
\end{tabular}
\par\smallskip\begin{minipage}{0.97\linewidth}\footnotesize Data: Open Source Asset Pricing release 2.0.0 \citep{chen2022open}, every predictor with a return in all 120 months of the window. Stationary bootstrap with 1{,}999 resamples shared by all tests; mean block length $q=5$ except where stated. Fixed: Hansen's studentization, the Politis--Romano standard deviation held fixed in every resample (at $q=10$, arch's default, this is the computation arch merged in October 2026). Bootstrap-$t$: each resample studentized by its own standard deviation. Totals over the windows: Fixed, $q=10$ 348, Fixed, $q=5$ 298, Bootstrap-$t$ 225, Reality Check 66.\end{minipage}
\end{table}



The simulations say which tests keep their level; they do not say whether the choice among them changes a conclusion on real data. As a small illustration, registered with v2b but without a prediction, we apply the tests to a classic search: \nz{emp.rules} moving-average timing rules on the US stock market \citep{brock1992simple,sullivan1999datasnooping}. The data are the daily market and risk-free returns in the Fama/French 3 Factors file of the Kenneth R. French Data Library \citep{fama1993common,french2026library}, from \nz{emp.first_day} to \nz{emp.last_day}; the file's SHA-256 is recorded with the results. Rule $L$, for $L=\nz{emp.rule_min},20,\dots,\nz{emp.rule_max}$ trading days, holds the market on day $t+1$ when a total-return index of the market closes above its $L$-day moving average on day $t$, and the risk-free asset otherwise. The series tested is each rule's return minus the market's, so the null hypothesis is that no rule beats buy-and-hold. We split the sample into \nz{emp.windows} non-overlapping windows of \nz{emp.window_len} trading days, the length of the simulated searches, from the first day on which every rule is defined. Each window is tested with the joint tests at mean block length 8 and \nz{emp.reps} resamples, with the fixed-variance and unstudentized tests also at block length 22, and with LET, the \v{S}id\'ak-adjusted $t$ test on the best rule's Sharpe ratio.

The tests agree on most windows: none rejects at 5\% in \nz{emp.none_reject} of the \nz{emp.windows}, and all reject in \nz{emp.all_reject} (Table~\ref{tab:empirical}; every window is in Table~\ref{tab:empirical-all}). In the other \nz{emp.split_windows} the verdict depends on the test, and the tests that the simulations found oversized reject most often. The fixed-variance SPA test rejects in \nz{emp.reject.spa_c_ge} windows, the Reality Check in \nz{emp.reject.rc}, the two re-studentized tests in \nz{emp.reject.spa_c_boot_t_ge} and \nz{emp.reject.spa_c_nb_ge}. The fixed-variance test rejects where the Reality Check does not in \nz{emp.only.spa_c_ge.not.rc} windows, \nz{emp.only.spa_c_ge.not.rc.list} (the two tests' $p$-values), and in \nz{emp.fixed_alone} of them neither re-studentized test rejects. At block length 22, arch's unstudentized computation rejects in \nz{emp.reject.spa_c_unstud_b22} windows against \nz{emp.reject.spa_c_unstud} at block length 8. LET rejects in only \nz{emp.reject.sidak_t}: its \v{S}id\'ak adjustment treats the twenty rules as independent, while rules with neighbouring lengths are almost the same rule, which makes it conservative, as in the correlated families of Section~\ref{sec:results}. Over the whole sample (\nz{emp.full.periods} days, mean block length \nz{emp.full.block}, and 22 for the two tests marked so), the \nz{emp.full.joint_count} joint tests' $p$-values range from \nz{emp.full.p.min_joint} to \nz{emp.full.p.max_joint}, at or just below the 5\% level, while LET's is \nz{emp.full.p.sidak_t}.

The illustration has the limits of its design. The rules are evaluated without transaction costs, so the series measure timing, not profit. The windows are not independent tests of one hypothesis: they are reported side by side, without a correction across windows, to show where the tests disagree, not to claim that any rule has skill. The one-sided test asks whether any rule beats buy-and-hold, and a rule that sits in the risk-free asset beats it in a falling market. Finally, the differences between tests come down to a handful of windows; the simulations, not this illustration, establish which test is calibrated.
